\subsection{\texttt{TSHR}}
\noindent\textbf{Bundle encoding.} UOP entry 61, opcode \texttt{0x1061}. Bundle: \texttt{B.OP + B.ARG + B.IOTI}. \texttt{B.OP.Opcode}=\texttt{0x1061}. \texttt{B.IOTI}: selector=\texttt{111 last}, \texttt{SrcTile0}=\texttt{src0}, \texttt{SrcTile1}=\texttt{src1}, \texttt{DstTile}=\texttt{dst}, \texttt{SizeCode}=profile tile size.\par\smallskip
{\small
\paragraph{PTO ISA description.}
Elementwise shift-right of two tiles.
\paragraph{Example.}
i8 tile [-8, 0, 8] shifted right by 1 (arithmetic) gives [-4, 0, 4]. ui8 tile [0b10000000, 0b01000000] shifted right by 1 gives [0b01000000, 0b00100000].
\par\smallskip
\paragraph{PTO ISA operation.}
For each element \texttt{(i, j)} in the valid region:

\[
\mathrm{dst}_{i,j} = \mathrm{src0}_{i,j} \gg \mathrm{src1}_{i,j}
\]
}\par\smallskip
\paragraph{Notes.}
The signed/unsigned interpretation follows the tile's declared element format. Shift amount >= element width produces zero (arithmetic shift of negative values produces all ones).
\paragraph{Microcode site table.}
{\scriptsize
\begin{longtable}{@{}R{0.055\linewidth}L{0.23\linewidth}L{0.31\linewidth}L{0.22\linewidth}@{}}
\toprule
Seq & Micro instruction & WO[63:32] fields & WM[31:0] fields \\
\midrule
0 & \texttt{u\_vec.vlds} & \texttt{opc=0x17, x=0, fl=mem, seq=0, kind=b} & \texttt{entry=61, cnt=2, res=vec, var=0} \\
1 & \texttt{u\_vec.vlds} & \texttt{opc=0x17, x=0, fl=mem, seq=1, kind=b} & \texttt{entry=61, cnt=2, res=vec, var=0} \\
2 & \texttt{u\_vec.vshr} & \texttt{opc=0x27, x=0, fl=alu, seq=2, kind=b} & \texttt{entry=61, cnt=2, res=vec, var=0} \\
3 & \texttt{u\_vec.vsts} & \texttt{opc=0x2A, x=0, fl=mem, seq=3, kind=b} & \texttt{entry=61, cnt=2, res=vec, var=0} \\
\bottomrule
\end{longtable}
}
\paragraph{Microcode body DSL.}
\begin{lstlisting}[style=ptocode]
def uop_tshr(src0, src1, dst):
    dtype = dst.element_type
    valid_rows, valid_cols = dst.valid_shape

    for row in range(0, valid_rows, 1):
        remained = valid_cols
        for col in range(0, valid_cols, pto.get_lanes(dtype)):
            mask, remained = pto.make_mask(dtype, remained)
            lhs = pto.vlds(src0[row, col:])
            rhs = pto.vlds(src1[row, col:])
            result = pto.vshr(lhs, rhs, mask)
            pto.vsts(result, dst[row, col:], mask)
    return
\end{lstlisting}
\Needspace{0.34\textheight}
